\Needspace{24\baselineskip}
\chapterrule
\section{逐个处理 Token，为什么仍看不到前文？}
\label{intro:21a-context-gap}
\chaptercompass{只逐行处理 Token，为什么无法理解前文？}{共享的逐位置函数只能看到本行输入，相同当前输入得到相同输出。}{$X:T\times d\to Y:T\times d'$，每个 $Y_t$ 只取决于 $X_t$}{逐位置变换的矩阵可学习，但不会自行产生跨行信息流}
\modelpipelinefocus{3}{上下文信息缺口}

独立处理每个 Token 的 MLP，只接收当前位置的向量。
如果下一 Token 的正确预测还取决于\textbf{前面说过什么}，
一个只看当前位置的函数能区分这些不同上下文吗？

\subsection{同一个当前词，两段不同的前文}

试着设计两条不同上下文：
\[
\text{“春天，花，开了”},\qquad
\text{“冬天，窗，开了”}.
\]
两段在最后位置都以“开了”结束，但后面的词
可以分别接不同对象或标点。这是说明“预测依赖上下文”的
\textbf{人为构造例子}，不是两句话的统计概率断言。

如果让一个相同的逐 Token MLP 只看“开了”的 Embedding，
它得到的输入都是同一个 $E[\text{开了}]$；
因为函数及参数也完全一样，两个上下文的输出就\textbf{必然一样}：
\[
f\bigl(E[\text{开了}]\bigr)
=f\bigl(E[\text{开了}]\bigr).
\]
它无法仅凭这一行区分不同前文。
要用 MLP 处理上下文，就要把前文也送进它的输入。

\subsection{朴素办法：把前几个 Token 直接拼起来}

可以把最近三个 Token 的向量首尾拼接成一条长输入，
交给一个 MLP。这样它确实能看到这三个位置，
但只允许固定长度时，碰到第四个、更早的第十个位置又怎么办？
如果不断扩展拼接长度，输入参数量和位置处理办法也随之改变。
固定窗口 MLP 能解决一些任务，仍没有自然地复用任意前文的机制。

\subsection{新的问题：怎样把历史信息送给当前位置？}

我们需要让后一个位置\textbf{能利用}前面位置的信息，
而不是只做独立行变换。
下一章先采用容易理解的方案：沿序列传递一份不断更新的状态。
这叫 RNN。再下一步才尝试不经过一长串状态，直接读取某个位置。

比较这三种机制时，要把任务长度、数据、参数量和计算代价一起列入评测。


\chapterchecktasks{21a-context-gap}{正文基础三题}{
\item \textbf{预测：}两句话的最后 Token 相同，但前文一个是否定、一个是肯定。只看最后一个 Token 的网络能否自动区分两句？为什么？\chapterinlinehint{只给当前 Token 时，两句无法区分。}
\item \textbf{手算：}两个序列末尾输入均为 $x=3$，同一个逐位置函数 $f(x)=2x+1$ 分别输出什么？\chapterinlinehint{两句都得到 $7$。}
\item \textbf{找错：}同学把过去三个 Token 拼起来当特征，就宣布“任意长历史都能无代价地读到”。还缺少哪两个工程条件？\chapterinlinehint{固定宽度拼接需限制长度或截断。}
}
\chapterrecallanswer{单位置函数不能从相同当前输入推断不同前文；上下文需要递推状态或跨位置读取机制。}



\begin{conclusion}
多维 Embedding 解决了“单个词怎样表示”，却不会自动提供跨位置上下文。
下一章试着让历史通过一份递推状态传到当前位置。
\end{conclusion}
